\documentclass[11pt]{article}
\usepackage[margin=1in]{geometry}
\usepackage{amsmath,amssymb,amsthm,mathtools}
\usepackage[T1]{fontenc}
\usepackage{lmodern,microtype,xurl}
\usepackage[colorlinks=true,linkcolor=blue,urlcolor=blue]{hyperref}
\hypersetup{pdftitle={A Wasserstein lower bound for local group-product networks},pdfauthor={ChatGPT-6 Astra, Kaarel Hänni, Dmitry Vaintrob}}
\newtheorem{theorem}{Theorem}[section]
\newtheorem{lemma}[theorem]{Lemma}
\newtheorem{corollary}[theorem]{Corollary}
\newtheorem{proposition}[theorem]{Proposition}
\theoremstyle{remark}\newtheorem{remark}[theorem]{Remark}
\newcommand{\E}{\mathbb E}
\newcommand{\R}{\mathbb R}
\newcommand{\norm}[1]{\left\lVert#1\right\rVert}
\newcommand{\PF}{P_{\mathrm F}}
\newcommand{\W}{W_2}
\newcommand{\cP}{\mathcal P}
\newcommand{\ent}{\mathcal H}
\title{A Wasserstein lower bound for local group-product networks}
\author{ChatGPT-6 Astra \and Kaarel H\"anni \and Dmitry Vaintrob}
\date{25 October 2026}
\begin{document}
\maketitle
\begin{abstract}
Consider a network computing a Boolean predicate of an ordered product of
$n$ letters in a fixed finite group. Suppose its physical layers merge
contiguous subwords through successive local modules, with disjoint data
channels between different modules in each layer band. Internal states may
retain the entire subword, have arbitrary dimension and precision, and use
arbitrary fixed-seed estimator features. No internal margin or product-state
encoding is assumed. We prove that the product of the layer Frobenius norms
is at least a constant times the exponential of one half the accumulated
entropy of the subword partitions. If every merge has arity at most $b$, this
is at least
$\exp((\log n)^2/(4\log b)+(\log n)/4)$, up to the stated input and margin
factor, for unbalanced as well as balanced hierarchies. The proof uses
Wasserstein distances between state distributions conditioned on the group
product. The lower bound concerns this physical architecture. It also applies
to local parity networks, although parity has a shallow polynomial-product
realization. It is therefore neither an intrinsic function-norm lower bound
nor a Bayesian sample lower bound.
\end{abstract}

\section{Purpose and the physical restriction}

A bounded-width branching program can multiply letters from a fixed finite
group in linear time. A balanced arithmetic formula can do the same with
polynomial size and logarithmic depth. These observations do not by themselves
supply a tanh network with small product of layer Frobenius norms: a physical
network must transmit its intermediate states through its actual matrices.
The issue is especially visible in local multiplication trees, which repeatedly
maintain many independent subword computations in parallel.

One possible response is to change the state code. A module might use a
high-dimensional code, keep auxiliary information about the original subword,
or store random-estimator values whose mean represents its group product.
The theorem below permits all such deterministic boundary states. Its
restriction is on dataflow: before their designated merge, different subwords
have disjoint data-dependent channels. Each merge band consists of at least
one actual matrix layer, and every layer that carries the states is charged.

This distinction matters. A free relabeling of two vectors as one concatenated
vector is not an additional physical layer. Nor may one factor a global linear
map virtually, assign the factors to an invented tree, and charge the resulting
logical depth. Conversely, arbitrary depth \emph{inside} a genuine local band
is permitted. The lower bound will not depend on that depth.

\subsection{Inputs and target}

Fix a finite group $G$ and $n\ge2$. The input letters $g_1,\ldots,g_n$
range independently over $G$. Letter $i$ is supplied through a fixed encoding
$E_i:G\to\R^{s_i}$. The principal case is an injective Boolean encoding,
$E_i(G)\subseteq\{-1,1\}^{s_i}$. A single augmented constant coordinate
may be provided in addition, as in the usual bias-free tanh architecture.

Fix a nonconstant predicate $f:G\to\{-1,1\}$, and put
\[
 p=\frac{|\{g\in G:f(g)=-1\}|}{|G|}\in(0,1).
\]
The network output $z$ is required to satisfy
\begin{equation}\label{eq:margin}
 f(g_1\cdots g_n)z(g_1,\ldots,g_n)\ge1
 \qquad\text{for every }(g_1,\ldots,g_n)\in G^n.
\end{equation}
The output magnitude may depend on the whole word. Only its required sign
depends on the group product. If the network is defined on a larger Boolean
cube, \eqref{eq:margin} is needed only on valid encoded words. In particular,
any full-cube extension satisfying the margin is covered.

Uniform letters are a device in the proof, used to define conditional state
distributions. The theorem is a deterministic representation statement, not
a learning guarantee for a specified sampling law.

\subsection{Partitions, bands, and arbitrary boundary codes}

Let
\[
 \cP_0,\cP_1,\ldots,\cP_T
\]
be successive partitions of $\{1,\ldots,n\}$ into contiguous intervals.
The first partition consists of singletons, the last consists of the whole
word, and each partition coarsens the preceding one. A block in $\cP_j$
is a parent; the consecutive blocks of $\cP_{j-1}$ that it contains are
its children. Unary carry modules are allowed.

At the cut associated with a block $v$, let $e_v(w_v)\in\R^{r_v}$ denote
its boundary code, where $w_v$ is the complete ordered subword in that block.
It is an arbitrary deterministic function of that subword. It may retain every
original letter and arbitrary auxiliary values. There is no assumption that
it is a function only of the product of the subword. The initial singleton
codes are the prescribed encodings $E_i$.

The bands, in order, partition all actual matrix layers preceding the final
readout. Each band is nonempty and consists of consecutive layers; no
physical layer is reused or omitted. Within band $j$, different parent modules have disjoint
data-dependent channels. The first matrix of each parent module receives
only the boundary codes of its children; subsequent matrices receive only
that module's data-dependent channels. Shared input-independent constants
are allowed. They may be treated as biases when deriving a lower bound,
although their actual weights are included in the network's norm.

All activations are coordinatewise $1$-Lipschitz. They may differ between
neurons and layers. Tanh is a special case; zero preservation is not needed
for the proof. The final readout is linear,
\[
 z=v^\top e_{\mathrm{root}}+c.
\]
Its constant can be supplied through ordinary charged constant channels.
Write
\[
 \PF=\prod_{\ell}\norm{W_\ell}_{\mathrm F}
\]
for the product over all actual weight matrices, including the readout.
Any additional charged constant weights only increase the lower estimates
used below. There are no uncharged input-dependent bypasses.

\section{The entropy-area theorem}

Give every block $u$ mass
\[
 \mu_u=|u|/n.
\]
At each cut these masses form a probability vector. Define its entropy in
natural units by
\[
 \ent(\cP_j)=-\sum_{u\in\cP_j}\mu_u\log\mu_u.
\]
For the input encodings, set
\[
 a_i^2=\E_{g\sim\mathrm{Unif}(G)}
       \norm{E_i(g)-\E E_i(g)}_2^2,
 \qquad A_0=\prod_{i=1}^n a_i^{1/n}.
\]
Whenever \eqref{eq:margin} is realizable in the stated architecture, all
$a_i$ are positive, as will follow from the proof. For Boolean $s_i$-bit
encodings, $a_i\le\sqrt{s_i}$.

\begin{theorem}[Entropy accumulated across physical bands]\label{thm:entropy}
Every network satisfying the preceding architecture assumptions and
\eqref{eq:margin} obeys
\begin{equation}\label{eq:entropy-bound}
 \PF\ge\frac{2\sqrt{p(1-p)}}{A_0}
   \exp\!\left(\frac12\sum_{j=1}^T\ent(\cP_{j-1})\right).
\end{equation}
Internal dimensions, code precision, retained history, and module depths
are unrestricted. Every term in the entropy sum corresponds to a nonempty
band of actual charged matrix layers.
\end{theorem}

The theorem separates two effects. Changes in the numerical scale of a
boundary code will telescope out. The entropy of the number and sizes of
simultaneously separate subword computations remains. The next corollary
makes this quantitative without requiring a balanced tree.

\begin{corollary}[Bounded merge arity]\label{cor:arity}
Suppose every parent has at most $b$ children, where $2\le b\le n$. Then
\begin{equation}\label{eq:arity-bound}
 \PF\ge\frac{2\sqrt{p(1-p)}}{A_0}
   \exp\!\left(\frac{(\log n)^2}{4\log b}+\frac{\log n}{4}\right).
\end{equation}
For a complete $b$-ary tree with $n=b^h$, this is
\[
 \PF\ge\frac{2\sqrt{p(1-p)}}{A_0}\,b^{h(h+1)/4}.
\]
\end{corollary}

For fixed arity the product is at least $\exp(\Omega((\log n)^2))$.
For arity of order $\log n$, the corresponding lower bound is
\[
 \exp\!\bigl(\Omega((\log n)^2/\log\log n)\bigr).
\]
When $b=n$ and there is only one band, the bound is only a constant times
$\sqrt n$. That case does allow an arbitrary global network inside the one
root band; no repeated local hierarchy is being imposed.

\section{Conditional Wasserstein signals}

For two finitely supported probability laws $\nu,\nu'$ on a Euclidean
space, their quadratic Wasserstein distance is
\[
 \W(\nu,\nu')^2
   =\min_{\pi\in\Pi(\nu,\nu')}
          \int\norm{x-y}_2^2\,d\pi(x,y),
\]
where $\Pi(\nu,\nu')$ is the set of couplings with the stated marginals.
The minimum exists. We use two elementary facts: applying an $L$-Lipschitz
map to both sides multiplies this squared distance by at most $L^2$;
and coupling mixtures with the same mixing variable bounds their squared
distance by the average squared distances of the coupled components.
Both facts follow by constructing the indicated coupling.

For a block $u$ and $g\in G$, let $\nu_{u,g}$ be the law of its boundary
code under uniform subwords conditioned on their product being $g$. Every
group product has positive probability, and the code laws have finite
support. Define
\begin{equation}\label{eq:signal}
 a_u^2=\frac1{2|G|^2}\sum_{g,h\in G}
                      \W(\nu_{u,g},\nu_{u,h})^2.
\end{equation}
The same notation agrees with the earlier input definition: at a leaf,
$\nu_{i,g}$ is the point mass at $E_i(g)$, and the average squared distance
of two independent points, divided by two, is their variance.

The conditional laws couple away arbitrary variation within a fixed product
class. No decoding map from code to product is needed. This is the reason the
argument permits complete histories and auxiliary estimator values.

\begin{lemma}[One-child contraction]\label{lem:contraction}
Suppose a parent module has code map
$F(e_1(w_1),\ldots,e_m(w_m))$. If its Lipschitz constant when only child
$i$'s code varies is at most $L_i$, then
\[
 a_{\mathrm{parent}}\le L_i a_i.
\]
\end{lemma}
\begin{proof}
We first give the binary case, varying the left child. Conditional on the
parent product being $a$, its complete right subword remains unconditionally
uniform. Indeed the independent left product is uniform on $G$, so each
fixed right subword contributes the same factor $1/|G|$ to the conditioning
probability. If the right subword has product $h$, the left product must be
$a h^{-1}$.

To couple parent code laws for products $a$ and $b$, choose the same uniform
right subword in both copies. Conditional on that word, couple the two left
codes optimally under
\[
 \nu_{\mathrm{left},a h^{-1}},\qquad
 \nu_{\mathrm{left},b h^{-1}}.
\]
Apply $F$ to these codes and the shared right code. This constructs a coupling
with the desired parent marginals and gives
\[
 \W(\nu_{\mathrm{parent},a},\nu_{\mathrm{parent},b})^2
 \le L_i^2\E_{w_{\mathrm{right}}}
       \W(\nu_{\mathrm{left},a h^{-1}},
           \nu_{\mathrm{left},b h^{-1}})^2.
\]
For every fixed $h$, right translation sends a uniform pair $(a,b)$ to a
uniform pair $(a h^{-1},b h^{-1})$. Average the last inequality over $a,b$
and use \eqref{eq:signal}.

For an arbitrary child of an $m$-ary parent, condition on all the other full
subwords and denote the products before and after that child by $L$ and $R$.
The other subwords again remain jointly uniform under parent-product
conditioning, because the missing child's product is independent uniform.
Its required product is $L^{-1}aR^{-1}$. The same construction and the
bijection $a\mapsto L^{-1}aR^{-1}$ prove the result.
\end{proof}

\begin{remark}
The parent code need not be constant on product classes, and its two
conditional laws need not have the same covariance or mean. The coupling
uses the full conditional laws. Replacing them by the means alone would
not justify the lemma for history-retaining states.
\end{remark}

\begin{lemma}[The output margin supplies a signal]\label{lem:root}
The root signal is positive, and
\[
 \norm v_2\ge\frac{2\sqrt{p(1-p)}}{a_{\mathrm{root}}}.
\]
\end{lemma}
\begin{proof}
For opposite labels $f(g)\ne f(h)$, every code in the support of
$\nu_{\mathrm{root},g}$ and every code in the support of
$\nu_{\mathrm{root},h}$ have readout values separated by at least two.
Their Euclidean separation is therefore at least $2/\norm v_2$. Every
coupling has at least this squared displacement, so
\[
 \W(\nu_{\mathrm{root},g},\nu_{\mathrm{root},h})^2
       \ge4/\norm v_2^2.
\]
The fraction of ordered opposite-label pairs is $2p(1-p)$. The factor
$1/2$ in \eqref{eq:signal} gives
$a_{\mathrm{root}}^2\ge4p(1-p)/\norm v_2^2$.
\end{proof}

Lemma~\ref{lem:contraction} propagates positivity from the root to every
preceding cut and every leaf. All signal ratios below are therefore legitimate.
The readout can alternatively be included in the last root band; then its
scalar output has signal at least $2\sqrt{p(1-p)}$.

\section{Physical norms and the entropy telescope}

Fix one band. For parent module $u$, partition its first variable-input
matrix into blocks $B_{u,i}$, one for each child. Let $T_u$ be the product of
the Frobenius norms of its remaining variable-channel matrices, with the
empty product equal to one. Coordinatewise activation contractions and the
operator-norm bound by Frobenius norm give a one-child Lipschitz constant
at most $T_u\norm{B_{u,i}}_{\mathrm F}$. Consequently
\begin{equation}\label{eq:module}
 P_u^2\ge T_u^2\sum_{i\text{ child of }u}
                   \norm{B_{u,i}}_{\mathrm F}^2
       \ge\sum_{i\text{ child of }u}(a_u/a_i)^2,
\end{equation}
where $P_u$ is the product of the module's variable-channel Frobenius norms.
Any constants omitted from this estimate can only increase actual matrix
norms. Shared input-independent channels may be frozen as biases when
bounding changes in the variable inputs.

At physical layer $t$ of the band, channel disjointness implies
\[
 \norm{W_t}_{\mathrm F}^2\ge\sum_u
                    \norm{B_{u,t}}_{\mathrm F}^2.
\]
Expanding a product of sums of nonnegative numbers gives
\begin{equation}\label{eq:band}
 P_{\mathrm{band}}^2
 \ge\prod_t\sum_u\norm{B_{u,t}}_{\mathrm F}^2
 \ge\sum_u\prod_t\norm{B_{u,t}}_{\mathrm F}^2
 =\sum_u P_u^2
 \ge\sum_{i\to u}(a_u/a_i)^2.
\end{equation}
This is why arbitrary additional physical depth inside a local band does not
weaken the conclusion. The estimate only requires the band to be nonempty
and all its variable-channel matrices to be charged.

\begin{proof}[Proof of Theorem~\ref{thm:entropy}]
Apply weighted arithmetic--geometric mean to the last sum in
\eqref{eq:band}, using weights $\mu_i$ on the child edges. These weights sum
to one because the child blocks form the pre-band partition. Thus
\begin{align*}
 \log P_{\mathrm{band}}
 &\ge\frac12\sum_{i\to u}\mu_i
                   \log\frac{(a_u/a_i)^2}{\mu_i}\\
 &=\sum_{i\to u}\mu_i(\log a_u-\log a_i)
           +\frac12\ent(\cP_{\mathrm{pre}}).
\end{align*}
For each parent, the child masses sum to its own mass. Hence, writing
$A_j=\prod_{u\in\cP_j}a_u^{\mu_u}$, the first term is exactly
$\log A_j-\log A_{j-1}$. Sum over bands. The result is
\[
 \prod_j P_{\mathrm{band},j}
 \ge\frac{a_{\mathrm{root}}}{A_0}
           \exp\!\left(\frac12\sum_j\ent(\cP_{j-1})\right).
\]
Multiply by the readout bound in Lemma~\ref{lem:root}. All internal signal
scales cancel, proving \eqref{eq:entropy-bound}.
\end{proof}

\begin{proof}[Proof of Corollary~\ref{cor:arity}]
The entropy chain rule for one coarsening gives
\[
 \ent(\cP_{j-1})-\ent(\cP_j)
 =\sum_{u\in\cP_j}\mu_u
       \ent\bigl((\mu_i/\mu_u)_{i\text{ child of }u}\bigr)
 \le\log b.
\]
Initial entropy is $A=\log n$ and final entropy is zero. With $B=\log b$
and $r=\lceil A/B\rceil$, the sum of pre-band entropies is at least
\[
 \sum_{j=0}^{r-1}(A-jB)
   =rA-\frac{Br(r-1)}2
   \ge\frac{A^2}{2B}+\frac A2.
\]
For the last inequality, set $r=A/B+\theta$, $0\le\theta<1$; the difference
is $B\theta(1-\theta)/2\ge0$. Substitute this estimate into
\eqref{eq:entropy-bound}. When $n=b^h$ and the tree is balanced, the cut
entropies are $h\log b,(h-1)\log b,\ldots,\log b$, giving the displayed
exact exponent.
\end{proof}

\begin{corollary}[Large interaction blocks must appear early]\label{cor:early}
Fix $G,f$ and bounded-size Boolean leaf encodings. Suppose networks in this
physical hierarchy satisfy $\PF\le C_0n^K$ for fixed $C_0>0$ and $K\ge0$. For every fixed
$0<\varepsilon<1$, a block containing at least $n^{1-\varepsilon}$ letters
must appear within $O((K+1)/\varepsilon)$ charged bands.
\end{corollary}
\begin{proof}
Theorem~\ref{thm:entropy} gives
$\sum_j\ent(\cP_{j-1})\le2K\log n+O(1)$. Thus only
$O((K+1)/\varepsilon)$ pre-band cuts can have entropy at least
$\varepsilon\log n$. If the hierarchy terminates earlier, its root already
has the required size. Otherwise some cut within that many bands has entropy
at most $\varepsilon\log n$. For any probability vector,
$\ent(\mu)\ge\log(1/\max\mu)$, so that cut has a block of mass at least
$n^{-\varepsilon}$.
\end{proof}

The conclusion concerns this hierarchy of physical interactions. It does not
assign an intrinsic interaction depth to an arbitrary fully connected network.

\section{Boolean examples and limits of the result}

\subsection{A fixed noncommutative group}

Take $G=A_5$, the even permutations of five objects, and define $f(g)=-1$
exactly when $g$ is a three-cycle. There are $20$ three-cycles among the
$60$ group elements, so $p=1/3$. Any injection of $A_5$ into six Boolean
bits gives $A_0\le\sqrt6$. Therefore every network for this predicate with
the stipulated maximum-arity-$b$ physical hierarchy satisfies
\[
 \PF\ge\frac{2}{3\sqrt3}
       \exp\!\left(\frac{(\log n)^2}{4\log b}+\frac{\log n}{4}\right).
\]
The claim holds for every Boolean extension on the unused four letter codes.
It does not prove that an unrestricted network for this word predicate must
have the displayed product.

The theorem rules out improving a local multiplication implementation merely
by distorting its group-state codes, keeping extra history, or placing more
hidden layers inside its local modules. Fixed-seed random-estimator features
are also deterministic subword codes and are covered when implemented through
the same charged dataflow. Ordinary arithmetic work is not being identified
with physical Frobenius product.

\subsection{Parity proves that the lower bound is architectural}

Take $G=\mathbb Z/2\mathbb Z$, encode each letter by one sign, and let $f$
be parity. Then $p=1/2$ and $A_0=1$. A binary local physical hierarchy must
have $\PF\ge2^{h(h+1)/4}$ for $n=2^h$. Yet parity has a shallow tanh
realization with polynomial product. Here is a direct construction.

For $x\in\{-1,1\}^n$, put $s(x)=\sum_i x_i$,
$\tau_j=-n+2j-1$, and $K=\log(8n)$. Use hidden neurons
\[
 t_j(x)=\tanh(K(s(x)-\tau_j)),\qquad j=1,\ldots,n,
\]
with all shifts supplied by the augmented constant input. Add a constant
neuron with value $c_0=\tanh1$. If $m$ coordinates are positive, then the
ideal signs $\eta_j$ equal $+1$ for $j\le m$ and $-1$ otherwise, and
\[
 1+\sum_{j=1}^n(-1)^j(1+\eta_j)=(-1)^m.
\]
Choose the linear readout
\[
 z(x)=2(-1)^n\left[1+\sum_{j=1}^n(-1)^j(1+t_j(x))\right],
\]
representing its constant term through $c_0$. Every preactivation before the
factor $K$ is an odd integer. Hence
$|t_j-\eta_j|\le2e^{-2K}$, and the output differs from
$2\prod_i x_i$ by at most $4ne^{-2K}<1$. Its signed margin exceeds one.
The first-layer Frobenius norm is $O(Kn^{3/2})$, including the shifts, and
the readout norm is $O(\sqrt n)$. Thus
\[
 \PF=O(n^2\log(n+1)).
\]
Its first hidden features access all input letters at once. It is a single
global merge, not the prescribed repeated local hierarchy, so there is no
contradiction.

\subsection{What remains outside the theorem}

Nonlocal interactions before a designated merge, globally mixed feature maps,
and other organizations of the network are not covered. Neither are
uncharged external rereads of input letters at later steps. If such values
are physically carried within child boundary states, they are covered as
retained history and their carrying matrices must be charged.

Finite ensembles are ordinary deterministic code vectors. The proof also
extends to deterministic Hilbert-valued boundary codes and charged
Hilbert--Schmidt maps whenever these objects define the physical realization;
the conditional distributions still have finite support over the finite input
set. It does not assert a result about arbitrary stochastic networks with
unspecified shared randomness.

The coupling lemma can similarly be formulated for finite operations that
are bijective in each argument and preserve uniform measure. No additional
application is needed here. The group case already shows the limitation of
local probabilistic branching-program implementations without introducing
internal robustness assumptions.

Finally, no prior distribution occurs in the proof. The Gaussian Bayesian
learner in the motivating problem is free to use fully connected global
features. An architecture-specific representation lower bound, even one
valid for arbitrary internal histories, does not establish a lower bound
for its posterior confidence or sample complexity.

\section*{Verification and provenance}

The conditional-law coupling, physical-band charging, entropy telescope,
and constants have conventional independent mathematical checks. This note
is not a Lean formalization. Its purpose is to retain a precise obstruction
to local compilation methods while leaving the unrestricted representation
and Bayesian questions open.

\end{document}
